\documentclass{article}
\usepackage{../math_template}

\author{Jonathan Kasongo}
\title{British Math Olympiad 2024 Round 1 --- P2/6}

\begin{document}
\maketitle

\begin{problem}{British Math Olympiad 2024 Round 1 --- P2/6}
A magician performs a trick with a deck of $n$ cards that are numbered from 1 to $n$. The
magician prepares for the trick by putting the cards in an order of her choosing. Then she
challenges a member of the audience to write an integer on a board. The magician turns over
the cards one by one, in their pre-arranged order. Every time the magician turns over a card,
the audience member multiplies the number on the board by $-1$, adds it to the number on
the card, writes the result on the board, and erases the old number. The magician guarantees
that, no matter which initial integer is chosen, the initial and final numbers will sum to 0.
Determine for which natural numbers $n$ the magician can perform the trick. You must both
prove that the trick is possible for the numbers you claim, and prove that it is not possible
for any other numbers.
\end{problem}

\begin{solution}{Solution}
The answer is that the trick is only possible if
$n \equiv 3 \ \text{(mod } 4)$. \\

Let's say that the initial number written on the board is $x$. Let's say
that the deck of cards is in the order $\sigma(1), \sigma(2), \ldots,
\sigma(n)$, where $\sigma$ is a permutation of $\{1,2,\ldots, n\}$. The
final number on the board is

\[
\begin{aligned}
&\phantom{=} \sigma(n) - \left[ \sigma(n-1) - \left[ \sigma(n-2) -
\cdots - \left[ \sigma(1) - x \right] \right] \cdots \right] \\
&= \sigma(n) - \sigma(n-1) + \cdots + (-1)^{n-1} \sigma(1) + (-1)^n x \\
\end{aligned}
\]

Note that if $n$ is even the final number will be of the form $a + x$, for
an integer $a$, and so the sum of the initial and final numbers is $a + 2x$
and so if $x \neq -\frac{a}{2}$ then that sum cannot be 0 and thus the trick
cannot be performed for any $x$. \\

So since $n$ is odd we just want to find $n$ such that there exists a
permutation $\sigma$ and

$$
\sigma(n) - \sigma(n-1) + \cdots + (-1)^{n-1} \sigma(1) = 0
$$

But since

$$
\sum_{r=0}^{n-1} (-1)^r \sigma(n-r) \equiv
\sum_{r=0}^{n-1} \sigma(n-r) \equiv
\frac{n(n+1)}{2} \ \text{(mod } 2)
$$

If $n \equiv 1 \ \text{(mod } 4)$ then $\frac{n(n+1)}{2}$ is odd and so
that sum can never be 0, since 0 is even. Thus we can only have
$n \equiv 3 \ \text{(mod } 4)$. Now to finish the problem we will exhibit
a construction that always works with such $n$. Consider the following
ordering of the deck of cards, we place each integer $1\leq k \leq n-1, \
k \neq \frac{n-1}{2}$ next to $k+1$, so that they will contribute $-1$
to the alternating sum, and then we stick the number $\frac{n-1}{2}$ on the
end of the list.

$$
1, 2, 3, 4, \ldots, \frac{n-1}{2} - 2, \frac{n-1}{2} - 1,  \frac{n-1}{2} + 1, \frac{n-1}{2} + 2, \ldots, n-1, n, \frac{n-1}{2}
$$

then the alternating sum of these numbers is
$$
\underbrace{\left[ (-1) + (-1) + \cdots (-1) \right]}_{\frac{n-1}{2} \ \text{copies}} + \frac{n-1}{2} = 0
$$

since there are $\frac{n-1}{2}$ pairs of integers $(k, k+1)$. $\blacksquare$
\end{solution}

\end{document}
